% Part: lambda-calculus
% Chapter: syntax
% Section: de-bruijn

\documentclass[../../../include/open-logic-section]{subfiles}

\begin{document}

\olfileid[id]{lam}{syn}{deb}

\olsection{Indeks de Bruijn}

Ekuivalensi-$\alpha$ sangatlah wajar, karena suku-suku yang
$\alpha$-ekuivalen ``bermakna sama.'' Bahkan, kita dapat memberikan sintaksis
bagi suku lambda yang tidak membedakan suku-suku yang dapat
dikonversi-$\alpha$ satu sama lain. Sintaksis yang paling dikenal mengganti
variabel dengan \emph{indeks de Bruijn}.

Ketika menulis $\lambd[x][M]$, kita menyatakan secara eksplisit bahwa $x$
adalah parameter fungsi, sehingga kita dapat menggunakan $x$ dalam $M$ untuk
merujuk kepada parameter tersebut. Akan tetapi, dalam indeks de Bruijn,
parameter tidak mempunyai nama dan acuan kepadanya dalam badan fungsi
dinyatakan oleh sebuah bilangan yang menunjukkan banyaknya tingkat abstraksi
antara acuan tersebut dan parameternya. Sebagai contoh, perhatikan
$\lambd[x][\lambd[y][y x]]$:
abstraksi luar mengikat variabel~$x$; abstraksi dalam mengikat variabel~$y$;
subsuku $y x$ berada dalam lingkup abstraksi dalam. Tidak ada abstraksi antara
$y$ dan abstraksinya, $\lambd[y]$, tetapi terdapat satu abstraksi antara $x$
dan abstraksinya, $\lambd[x]$. Jadi, kita menulis $0\, 1$ untuk $y x$, dan
$\lambd[][\lambd[][01]]$ untuk seluruh suku tersebut.

\begin{defn}
  Suku de Bruijn didefinisikan secara induktif sebagai berikut:
  \begin{enumerate}
  \item $n$, dengan $n$ sebarang bilangan asli.
  \item $PQ$, dengan $P$ dan $Q$ keduanya suku de Bruijn.
  \item $\lambd[][N]$, dengan $N$ suatu suku de Bruijn.
  \end{enumerate}
\end{defn}

Penerjemahan formal dari suku lambda biasa ke suku berindeks de Bruijn adalah
sebagai berikut:
\begin{defn}
  \begin{align*}
    F_\Gamma(x) &= \Gamma(x) \\
    F_\Gamma(PQ) &= F_\Gamma(P)F_\Gamma(Q) \\
    F_\Gamma(\lambd[x][N]) &= \lambd[][F_{x,\Gamma}(N)]
  \end{align*}
  dengan $\Gamma$ suatu daftar variabel yang diindeks mulai dari nol, dan
  $\Gamma(x)$ menyatakan posisi variabel $x$ dalam $\Gamma$. Sebagai contoh,
  jika $\Gamma$ adalah $x,y,z$, maka $\Gamma(x)$ adalah $0$ dan $\Gamma(z)$
  adalah $2$.

  $x,\Gamma$ menyatakan daftar yang dihasilkan dengan menambahkan $x$ ke awal
  $\Gamma$; misalnya, dengan melanjutkan $\Gamma$ dalam contoh terakhir,
  $w,\Gamma$ adalah $w,x,y,z$.
\end{defn}

Pemulihan suku lambda standar dari suku de Bruijn dilakukan sebagai berikut:

\begin{defn}
  \begin{align*}
    G_\Gamma(n) &= \Gamma[n] \\
    G_\Gamma(PQ) &= G_\Gamma(P) G_\Gamma(Q) \\
    G_\Gamma(\lambd[][N]) &= \lambd[x][G_{x,\Gamma}(N)]
  \end{align*}
  dengan $\Gamma$ sekali lagi suatu daftar variabel yang diindeks mulai dari
  nol, dan $\Gamma[n]$ menyatakan variabel pada posisi~$n$. Sebagai contoh,
  jika $\Gamma$ adalah $x,y,z$, maka $\Gamma[1]$ adalah $y$.

  Variabel $x$ dalam persamaan terakhir dipilih sebagai sebarang variabel
  yang tidak terdapat dalam $\Gamma$.
\end{defn}

Di sini kita memberikan beberapa hasil tanpa pembuktian:

\begin{prop}
  Jika $M \aconvone M'$, dan $\Gamma$ adalah sebarang daftar yang memuat
  $\FV{M}$, maka $F_\Gamma(M) \eqs F_\Gamma(M')$.
\end{prop}

\end{document}
